Small transformers trained on algorithmic tables can first memorise the training set and only much later generalise, a delay known as grokking~\cite{power2022grokking}. Modular addition is the standard testbed: a one-layer transformer ends up using Fourier features of the residues~\cite{nanda2023progress}, simple MLPs show the same jump~\cite{gromov2023grokking}, and the time to generalise depends strongly on the weight decay and the fraction of the table that is observed~\cite{power2022grokking,liu2022understanding}. Because grokking is slow, any cheap intervention that shortens the delay is of practical interest.

Curriculum learning~\cite{bengio2009curriculum,wang2020survey} is a natural candidate: present easy examples first. For modular addition the notion of ``easy'' is not obvious, since the operation is a group law that is symmetric in the residues, but operand size is a natural proxy that a practitioner would try. The question of this paper is therefore narrow and empirical: does ordering the training pairs by operand size (a curriculum), or by decreasing operand size (an anti-curriculum), change the number of steps a one-layer transformer needs to reach 95\% held-out accuracy compared with uniform sampling, and how does this interact with weight decay and training fraction?

Contributions, all at CPU scale: (i) a controlled comparison of the three samplers with identical model, optimiser, budget and splits (Sections~\ref{sec:method}--\ref{sec:setup}); (ii) a weight-decay $\times$ training-fraction grid and a pacing/shuffled-order sweep that probe the interaction and the mechanism (Section~\ref{sec:ablations}); (iii) a mostly negative result: the curriculum was not faster than uniform sampling in our setting, with the evidence limited by small seed counts and by many runs not generalising within the budget. We state where grokking was and was not reached.
